\subsection{辅助生殖期间的性生活：常见疑问速答}

体外受精（IVF）、人工授精（IUI）等技术让许多家庭迎来新生命，也让另一个问题浮出水面：\textbf{治疗期间到底还能不能同房？}性健康视角下，本节回答几个高频疑问。各阶段的完整时序指引见女性卷"生育问题与辅助生殖技术"一章；此处补充本书（综合卷）读者的常见困惑。

\vspace{0.5em}
\noindent\textbf{高频疑问三步答}：

\begin{itemize}
	\item \textbf{"促排期间能同房吗？"}——可以，但\textbf{必须严格避孕}。促排药物若在意外妊娠中继续暴露，可能干扰早期胚胎发育；卵巢因多卵泡发育而增大，此阶段也应避免深插入体位下的腹部受压与姿势骤变（警惕卵巢扭转）。
	\item \textbf{"移植后能同房吗？会不会把胚胎'颠'出来？"}——\textbf{没有证据表明移植后的性生活会导致胚胎着床失败或流产}，但多数生殖中心仍建议移植后至验孕前暂停，主要目的是减少出血诱因、避免"万一失败是不是怪那次"的心理负担。
	\item \textbf{"取卵前后要注意什么？"}——取卵前 1--2 天起暂停同房（降低穿刺区感染与出血风险），取卵后通常建议 1--2 周内避免同房、盆浴与游泳，待穿刺针道闭合；出现阴道出血增多、腹痛加剧或发热应及时联系中心。
\end{itemize}

\vspace{0.5em}
\noindent\textbf{治疗期的亲密不应只有"插入"}：两周等待期、术后恢复期里，拥抱、依偎、互相按摩、共同期待，都是维系联结的方式。不少夫妇事后回忆，这段"只许温柔"的时期反而成了难得的慢时光。

\vspace{0.5em}
\noindent\textbf{把生育与性暂时分开}：当每次同房都承载"必须怀上"的重量，性唤起与高潮都会被焦虑抑制。治疗周期内，如条件允许（医嘱未禁止），保留一些"只为亲密、不为生育"的接触，有助于保护关系本身。若经历失败周期后持续回避性、影响关系，建议寻求心理支持。

\begin{tcolorbox}[colback=pink!5!white,colframe=red!50!black,title=贴心小叮咛]
辅助生殖期间的同房问题，是生殖中心每天都会被问到的\textbf{常规问题}——不必"不好意思开口"。任何阶段拿不准，直接问主管医生，比在网上搜答案可靠得多。详见女性卷"辅助生殖各阶段的性生活指引"一节。
\end{tcolorbox}
